%% FORMLÆRE: MATHEMATICA
%% The Formal Mathematical Companion to Tractatus Formae
%% Iver Raknes Finne, Institute of Design, AHO
%%
%% This document contains the complete axiomatic-deductive structure
%% underlying Tractatus Formae: all definitions, postulates, theorems,
%% and proofs referenced via [M: ...] in the tractate.
%%
%% Compiled with: pdflatex → biber → pdflatex × 2

\documentclass[11pt, a4paper]{article}

%% ── Encoding & fonts ──────────────────────────────────────────────
\usepackage[T1]{fontenc}
\usepackage[utf8]{inputenc}
\usepackage{ebgaramond}            % Match Tractatus Formae typography

%% ── Mathematics ───────────────────────────────────────────────────
\usepackage{amsmath, amssymb, amsthm, mathtools}
\usepackage{thmtools}
\usepackage{bm}                    % Bold math
\usepackage{dsfont}                % Double-stroke font for \mathds

%% ── Layout ────────────────────────────────────────────────────────
\usepackage[
  top=30mm, bottom=35mm,
  left=32mm, right=32mm
]{geometry}
\usepackage{setspace}
\onehalfspacing
\usepackage{titlesec}
\usepackage{enumitem}
\usepackage{booktabs}
\usepackage{longtable}

%% ── References & links ───────────────────────────────────────────
\usepackage[
  backend=biber,
  style=authoryear-comp,
  sorting=nyt,
  maxcitenames=2
]{biblatex}
\addbibresource{../references.bib}
\usepackage[hidelinks]{hyperref}
\usepackage{cleveref}

%% ── Theorem environments ─────────────────────────────────────────
\theoremstyle{definition}
\newtheorem{definition}{Definition}[section]
\newtheorem{postulate}[definition]{Postulate}
\newtheorem{proposition}[definition]{Proposition}

\theoremstyle{plain}
\newtheorem{theorem}[definition]{Theorem}
\newtheorem{lemma}[definition]{Lemma}
\newtheorem{corollary}[definition]{Corollary}

\theoremstyle{remark}
\newtheorem{remark}[definition]{Remark}
\newtheorem{notation}[definition]{Notation}
\newtheorem{example}[definition]{Example}

%% ── Custom commands ──────────────────────────────────────────────
\newcommand{\MC}{M^{\mathsf{I}}_C}          % Morphospace
\providecommand{\Rn}{}
\renewcommand{\Rn}{\mathbf{R}^n}            % n-dim Euclidean space
\providecommand{\R}{}
\renewcommand{\R}{\mathbf{R}}               % Reals
\newcommand{\NN}{\mathbf{N}}                 % Naturals
\newcommand{\Pareto}{\mathcal{P}}            % Pareto front
\newcommand{\FitLand}{\mathcal{L}}           % Fitness landscape
\newcommand{\Nav}{\mathcal{N}}               % Navigator
\newcommand{\CogBound}{\partial\!\Nav}       % Cognitive boundary
\newcommand{\inner}[2]{\langle #1, #2 \rangle}
\newcommand{\norm}[1]{\lVert #1 \rVert}
\newcommand{\abs}[1]{\lvert #1 \rvert}
\DeclareMathOperator{\grad}{\nabla}
\DeclareMathOperator{\supp}{supp}
\DeclareMathOperator{\Int}{int}
\DeclareMathOperator{\MI}{I}                 % Mutual information
\DeclareMathOperator{\Ent}{H}                % Shannon entropy
\DeclareMathOperator{\argmax}{arg\,max}

%% ── Section formatting (small caps, centered) ────────────────────
\titleformat{\section}
  {\normalfont\large\scshape\centering}
  {\thesection}{1em}{}
\titleformat{\subsection}
  {\normalfont\normalsize\itshape}
  {\thesubsection}{1em}{}

%% ── Tractatus cross-reference command ────────────────────────────
\newcommand{\TF}[1]{\textnormal{[TF\,#1]}}  % Cross-ref to Tractatus Formae

%% ══════════════════════════════════════════════════════════════════
\begin{document}

%% ── Title page ───────────────────────────────────────────────────
\begin{titlepage}
\centering
\vspace*{4cm}
{\Large\scshape Forml\ae re: Mathematica}

\vspace{1cm}
{\large\itshape The Axiomatic--Deductive Companion to\\[4pt]
Tractatus Formae}

\vspace{2cm}
{\normalsize Iver Raknes Finne}

\vspace{0.5cm}
{\small Institute of Design, Oslo School of Architecture and Design}

\vspace{4cm}
{\small 2026}

\vfill
{\footnotesize
Cross-references to the tractate are marked \TF{\ldots}.\\
Every definition, postulate, and theorem listed in the Tractatus Formae\\
appendix ``Formell korrespondanse'' finds its proof here.}
\end{titlepage}

%% ── Table of contents ────────────────────────────────────────────
\tableofcontents
\newpage

%% ══════════════════════════════════════════════════════════════════
%% CHAPTER 0: PREAMBLE
%% ══════════════════════════════════════════════════════════════════
\section*{Preamble}
\addcontentsline{toc}{section}{Preamble}

This document provides the complete formal apparatus underlying
\textit{Tractatus Formae}. Where the tractate argues in natural language,
this companion argues in set theory, topology, optimization theory,
information theory, and dynamical systems. Every definition, postulate,
and theorem referenced via \textup{[M:\,\ldots]} in the tractate appears
here with its full proof.

\medskip\noindent\textbf{Logical architecture.}
The theory rests on five \emph{postulates}---empirically falsifiable
assumptions about the world---from which all theorems follow deductively.
The postulates are:
\begin{enumerate}[nosep]
  \item \textbf{Postulate~2.3}\; (Multiple independent pressures)
  \item \textbf{Postulate~2.6}\; (Conflicting gradients)
  \item \textbf{Postulate~4.1}\; (Dynamic landscape)
  \item \textbf{Postulate~5.1}\; (Geometric material signature)
  \item \textbf{Postulate~7.1}\; (Substrate-independent navigation)
\end{enumerate}
Two additional postulates appear in the variational extension:
Postulate~9.1 (path-dependent fitness) and its empirical derivative
(spreading hypothesis, Theorem~9.4).

\medskip\noindent\textbf{Notation.}
We write $\MC$ for the morphospace of class $C$ (Def.~1.1),
$s_i$ for the $i$-th selection pressure (Def.~2.1),
$\bm{s}$ for the selection-pressure vector (Def.~2.2),
$f$ for the fitness function (Def.~3.1),
$\FitLand$ for the fitness landscape (Def.~3.2),
and $\Nav$ for a navigator (Def.~6.1).
The gradient of a scalar field $g$ is $\grad g$.
Shannon entropy is $\Ent$, mutual information is $\MI$.

\newpage
%% ══════════════════════════════════════════════════════════════════
%% CHAPTER 1: MORPHOSPACE
%% ══════════════════════════════════════════════════════════════════
\section{Morphospace}\label{sec:morphospace}

\begin{definition}[Morphospace]\label{def:1.1}
Let $C$ be a functional class of objects (e.g.\ chairs, knives, vessels)
and let $x_1, x_2, \ldots, x_n$ be $n$ measurable geometric properties
that characterize objects in $C$ (height, width, mass, curvature,
symmetry, compactness, etc.).

For any object $a \in C$, its \emph{attribute vector} is
\[
  \bm{x}(a) = \bigl(x_1(a),\; x_2(a),\; \ldots,\; x_n(a)\bigr) \in \Rn.
\]
The \emph{morphospace} of class $C$ is the set
\[
  \MC \;=\; \bigl\{\, \bm{x} \in \Rn
    \;\big|\; \bm{x} \text{ satisfies the physical and functional
    constraints of } C \,\bigr\}
\]
equipped with the Euclidean metric $d(\bm{x}, \bm{y}) = \norm{\bm{x} - \bm{y}}_2$.
\end{definition}

\begin{definition}[Feasibility constraints]\label{def:1.2}
Let $\{g_j\}_{j=1}^{m}$ be a finite set of constraint functions
$g_j : \Rn \to \R$ encoding physical laws, material limits, and
functional requirements. The morphospace is the feasible set:
\[
  \MC \;=\; \bigl\{\, \bm{x} \in \Rn
    \;\big|\; g_j(\bm{x}) \leq 0 \;\;\forall\, j = 1, \ldots, m
  \,\bigr\} \;\subseteq\; \Rn.
\]
We assume each $g_j$ is continuous, so that $\MC$ is a closed subset
of\/ $\Rn$. If additionally the property ranges are bounded (which holds
for all physical objects), then $\MC$ is \emph{compact}.
\end{definition}

\begin{remark}
Compactness of\/ $\MC$ is not a mathematical convenience but a physical
fact: no chair has infinite height, negative mass, or unbounded curvature.
Compactness guarantees that continuous functions on $\MC$ attain their
extrema (Weierstrass), which is essential for the existence of optima in
the fitness landscape (\S\ref{sec:landscape}).
\end{remark}

\begin{definition}[Region typology]\label{def:1.3}
The full ambient space $\Rn$ partitions into three regions relative to
the class $C$:
\begin{enumerate}[label=(\roman*),nosep]
  \item \textbf{Busett (occupied) region} $B$:
    \[
      B \;=\; \bigl\{\, \bm{x} \in \MC
        \;\big|\; \bm{x} = \bm{x}(a) \text{ for at least one
        realized object } a \in C \,\bigr\}.
    \]
    Formally, $B = \supp(\mu)$ where $\mu$ is the empirical measure
    on $\MC$ induced by the population of realized objects.

  \item \textbf{Open region} $O$:
    \[
      O \;=\; \MC \setminus \overline{B}
    \]
    Positions that are physically feasible but historically unrealized.

  \item \textbf{Forbidden region} $F$:
    \[
      F \;=\; \Rn \setminus \MC
    \]
    Positions that violate at least one constraint $g_j$.
\end{enumerate}
Thus $\Rn = B \cup O \cup F$ (up to boundary effects), and
$\MC = B \cup O$.
\end{definition}

\begin{proposition}[Explanation requirement]\label{prop:1.4}
If $\abs{B} > 1$ (the occupied region contains more than one distinct
position), then the fact that a particular object $a$ occupies position
$\bm{x}(a)$ and not some other feasible position $\bm{y} \in \MC$
requires explanation.

\begin{proof}
Since $\MC$ is a compact subset of $\Rn$ with $n \geq 1$ and $\abs{B} > 1$,
there exist at least two distinct feasible positions. By the principle of
sufficient reason (Leibniz): if there were no cause favouring $\bm{x}(a)$
over $\bm{y}$, the occupation of $\bm{x}(a)$ would be arbitrary.
But the empirical distribution over $B$ is non-uniform
(Postulate~\ref{post:2.3}), so the explanation must come from
\emph{selection pressures} that favour certain positions. \qedhere
\end{proof}
\end{proposition}

\TF{1.1, 1.11, 1.12, 1.13, 1.14}

\newpage
%% ══════════════════════════════════════════════════════════════════
%% CHAPTER 2: SELECTION PRESSURES
%% ══════════════════════════════════════════════════════════════════
\section{Selection Pressures}\label{sec:pressures}

\begin{definition}[Selection pressure]\label{def:2.1}
A \emph{selection pressure} on morphospace $\MC$ is a continuous
function
\[
  s_i : \MC \longrightarrow \R
\]
that assigns to each position $\bm{x}$ a real value $s_i(\bm{x})$
expressing how well $\bm{x}$ satisfies the $i$-th criterion.
Higher values indicate better satisfaction. The gradient
$\grad s_i(\bm{x})$, where it exists, points in the direction of
steepest improvement of the $i$-th criterion.
\end{definition}

\begin{remark}
A selection pressure is not a rule (``do this'') but a \emph{gradient}
(``this direction is better''). It is continuous, not discrete;
directional, not prescriptive. It produces a \emph{tendency}---a
statistical overrepresentation of certain positions---not a deterministic
outcome. \TF{2.11, 2.12}
\end{remark}

\begin{definition}[Selection-pressure vector]\label{def:2.2}
Given $k$ selection pressures $s_1, \ldots, s_k$ acting on $\MC$,
the \emph{selection-pressure vector} is the map
\[
  \bm{s} : \MC \longrightarrow \R^k, \qquad
  \bm{s}(\bm{x}) = \bigl(s_1(\bm{x}),\; s_2(\bm{x}),\; \ldots,\;
  s_k(\bm{x})\bigr).
\]
Each component $s_i$ represents an independent criterion. \TF{2.14}
\end{definition}

\begin{postulate}[Multiple independent pressures]\label{post:2.3}
For every functional class $C$ with morphospace $\MC$:
\begin{enumerate}[label=(\alph*),nosep]
  \item $k \geq 2$: at least two selection pressures act on $\MC$.
  \item There exist at least two pressures $s_i, s_j$ that are
    \emph{statistically independent} over $\MC$: knowing $s_i(\bm{x})$
    does not determine $s_j(\bm{x})$.
\end{enumerate}

\medskip\noindent
\textit{Falsification criterion.}
The postulate fails if there exists a class $C$ in which a single
selection pressure $s_1$ is sufficient to explain all observed variation
in $B$---i.e., $B$ is one-dimensional in the image of $\bm{s}$ and no
second independent pressure can be identified. \TF{2.2}
\end{postulate}

\begin{theorem}[Variation under constant function]\label{thm:2.4}
Let $C$ be a functional class with $k \geq 2$ selection pressures.
Let $s^f$ denote the functional pressure (how well the object serves its
purpose). Since all objects in $C$ share the same function,
$s^f(\bm{x})$ is constant for all $\bm{x} \in \MC$ that satisfy the
functional requirement. Then:
\begin{enumerate}[label=(\roman*),nosep]
  \item $s^f$ cannot explain any variation in form within $C$.
  \item All observed form variation is attributable to the remaining
    pressures $s^k$ ($k \neq f$).
\end{enumerate}

\begin{proof}
\textit{(i)}\; By definition of a functional class, every $\bm{x} \in B$
satisfies the functional criterion. Thus $s^f$ is constant on $B$:
$s^f(\bm{x}) = c$ for all $\bm{x} \in B$. A constant function has
zero gradient: $\grad s^f = \bm{0}$ everywhere on $B$. Therefore $s^f$
induces no directional preference and cannot explain why any position is
occupied rather than another.

\textit{(ii)}\; By Postulate~\ref{post:2.3}, $k \geq 2$ and at least one
pressure $s^j$ ($j \neq f$) is independent of $s^f$. Since $s^f$ is
constant, the observed variation in $B$ (which is non-trivial since
$\abs{B} > 1$) must be explained by the gradients of the remaining
pressures $s^j$, $j \neq f$. \qedhere
\end{proof}
\end{theorem}

\TF{2.21}

\begin{definition}[Functional delimitation]\label{def:2.5}
Let $F_C : \Rn \to \{0, 1\}$ be the indicator function of the
functional requirement:
\[
  F_C(\bm{x}) =
  \begin{cases}
    1 & \text{if } \bm{x} \text{ satisfies the functional requirement
        of class } C, \\
    0 & \text{otherwise.}
  \end{cases}
\]
Then $\MC = \{\bm{x} \in \Rn : F_C(\bm{x}) = 1 \text{ and }
g_j(\bm{x}) \leq 0 \; \forall j\}$.
The function defines the class and thereby the morphospace.
Within $\MC$, $F_C \equiv 1$, so the function is \emph{silent}:
it delimits but does not differentiate. \TF{2.4}
\end{definition}

\begin{postulate}[Conflicting gradients]\label{post:2.6}
For at least one pair $(s_i, s_j)$ of selection pressures, there exists
a region $R \subseteq \MC$ with $\abs{R} > 0$ (positive measure) such
that
\[
  \inner{\grad s_i(\bm{x})}{\grad s_j(\bm{x})} < 0
  \qquad \forall\, \bm{x} \in R.
\]
That is: improving one property worsens another.

\medskip\noindent
\textit{Falsification criterion.}
The postulate fails if all gradient pairs are positively correlated
everywhere in $\MC$---i.e., all pressures can be simultaneously improved
at every point. \TF{2.3}
\end{postulate}

\begin{theorem}[Compromise]\label{thm:2.7}
Under Postulates~\ref{post:2.3} and \ref{post:2.6}, there exists no
position $\bm{x}^* \in \MC$ at which all selection pressures are
simultaneously maximized. Every realized position $\bm{x} \in B$
is a \emph{compromise}.

\begin{proof}
Suppose for contradiction that $\bm{x}^*$ is a simultaneous maximum:
$\grad s_i(\bm{x}^*) = \bm{0}$ for all $i = 1, \ldots, k$ and each
$s_i$ has a local maximum at $\bm{x}^*$.

By Postulate~\ref{post:2.6}, there exists a region $R$ with positive
measure where $\inner{\grad s_i}{\grad s_j} < 0$ for some pair $(i, j)$.
At $\bm{x}^*$, if all gradients vanish, then the inner product is zero,
not negative. So $\bm{x}^*$ cannot lie in the interior of $R$.

More strongly: in the region $R$, any direction $\bm{v}$ that increases
$s_i$ (i.e., $\inner{\grad s_i}{\bm{v}} > 0$) necessarily decreases
$s_j$ (since $\grad s_j$ has a negative component along $\grad s_i$).
Therefore no critical point in $R$ can be a maximum for both $s_i$ and
$s_j$ simultaneously.

Since the feasible set $\MC$ is compact and each $s_i$ is continuous,
each $s_i$ attains its individual maximum on $\MC$ (Weierstrass). But
these maxima occur at different points:
$\bm{x}_i^* = \argmax_{\bm{x} \in \MC} s_i(\bm{x})$ with
$\bm{x}_i^* \neq \bm{x}_j^*$ in general, since the gradients conflict.
No single point maximizes all simultaneously. \qedhere
\end{proof}
\end{theorem}

\TF{2.31}

\begin{theorem}[Multiple valid compromises]\label{thm:2.8}
Under the conditions of Theorem~\ref{thm:2.7}, and given that
$\abs{B} > 1$ (Proposition~\ref{prop:1.4}), there exist multiple
distinct compromises: different weightings of the pressures produce
different optima.

\begin{proof}
Let $\Phi_{\bm{w}} : \R^k \to \R$ be a monotonically increasing
aggregation function parameterized by weight vector
$\bm{w} = (w_1, \ldots, w_k) \in \Delta^{k-1}$ (the $(k{-}1)$-simplex).
Define the $\bm{w}$-weighted fitness:
\[
  f_{\bm{w}}(\bm{x}) = \sum_{i=1}^{k} w_i \, s_i(\bm{x}).
\]
Since $\MC$ is compact and $f_{\bm{w}}$ is continuous,
$f_{\bm{w}}$ attains its maximum at some
$\bm{x}^*(\bm{w}) \in \MC$.

For distinct weight vectors $\bm{w} \neq \bm{w}'$, the maxima
generically differ: $\bm{x}^*(\bm{w}) \neq \bm{x}^*(\bm{w}')$.
To see this, note that at $\bm{x}^*(\bm{w})$, the first-order
condition requires
\[
  \sum_{i=1}^k w_i \, \grad s_i(\bm{x}^*) = \bm{0}.
\]
Under Postulate~\ref{post:2.6}, the gradients $\grad s_i$ are not
positively proportional, so different $\bm{w}$ generically yield
different critical points. Therefore the set of compromises
$\{\bm{x}^*(\bm{w}) : \bm{w} \in \Delta^{k-1}\}$ contains more than
one element, and this set traces out the \emph{Pareto front}
$\Pareto \subseteq \MC$.

The difference between two compromises under the same function is not a
defect but reflects different balancings of the same forces. \qedhere
\end{proof}
\end{theorem}

\TF{2.32}

\newpage
%% ══════════════════════════════════════════════════════════════════
%% CHAPTER 3: FITNESS LANDSCAPE
%% ══════════════════════════════════════════════════════════════════
\section{Fitness Landscape}\label{sec:landscape}

\begin{definition}[Fitness function]\label{def:3.1}
The \emph{fitness function} (or \emph{adaptation function}) is a
continuous map
\[
  f : \MC \longrightarrow \R, \qquad
  f(\bm{x}) = \Phi\bigl(s_1(\bm{x}),\; s_2(\bm{x}),\; \ldots,\;
  s_k(\bm{x})\bigr)
\]
where $\Phi : \R^k \to \R$ is a smooth aggregation function,
monotonically non-decreasing in each argument. The choice of $\Phi$
is not unique; different aggregations encode different relative weightings
of the pressures. \TF{3.1}
\end{definition}

\begin{remark}
In the simplest model $\Phi$ is linear:
$f(\bm{x}) = \sum_i w_i \, s_i(\bm{x})$ with $w_i \geq 0$,
$\sum w_i = 1$. But $\Phi$ may be nonlinear if pressures interact
epistatically---e.g., if the effectiveness of ergonomic optimization
depends on the material chosen. \TF{3.11}
\end{remark}

\begin{definition}[Fitness landscape]\label{def:3.2}
The \emph{fitness landscape} is the graph of $f$ over $\MC$:
\[
  \FitLand \;=\; \bigl\{\, (\bm{x},\, f(\bm{x}))
    \;\big|\; \bm{x} \in \MC \,\bigr\}
  \;\subseteq\; \R^{n+1}.
\]
Local maxima of $f$ are \emph{peaks} (hills). Local minima are
\emph{valleys}. Saddle points are \emph{ridges} separating peaks. \TF{3.2}
\end{definition}

\begin{theorem}[Multiple peaks]\label{thm:3.3}
Under Postulates~\ref{post:2.3} and \ref{post:2.6}, the fitness function
$f$ generically has at least two local maxima on $\MC$.

\begin{proof}
We establish this via two independent arguments.

\medskip\noindent\textit{Argument 1 (Morse-theoretic).}\;
Assume $f$ is a Morse function on $\MC$ (all critical points are
non-degenerate; this holds generically by the Morse--Sard theorem).
Let $m_\lambda$ denote the number of critical points of index $\lambda$
(eigenvalues of the Hessian with negative sign). By the Morse
inequalities on a compact manifold with boundary:
\[
  m_0 - m_1 + m_2 - \cdots \;=\; \chi(\MC)
\]
where $\chi(\MC)$ is the Euler characteristic. For a compact convex
body in $\Rn$, $\chi = 1$, so $m_0 \geq 1$ (at least one local
minimum) and $m_n \geq 1$ (at least one local maximum).

By Postulate~\ref{post:2.6}, the Hessian of $f$ at any critical point
reflects conflicting curvatures from different pressures. In the region
$R$ where $\inner{\grad s_i}{\grad s_j} < 0$, the aggregated landscape
develops saddle points separating basins. By Morse theory, the existence
of a saddle point (index 1) between two regions implies
$m_n \geq m_1 + 1 \geq 2$ local maxima.

\medskip\noindent\textit{Argument 2 (Direct construction).}\;
Let $s_i, s_j$ have conflicting gradients in $R$ (Postulate~\ref{post:2.6}).
Consider the one-parameter family of fitness functions
$f_t(\bm{x}) = t \, s_i(\bm{x}) + (1-t) \, s_j(\bm{x})$ for
$t \in [0,1]$. At $t = 1$, $f_1 = s_i$ has its maximum at
$\bm{x}_i^* = \argmax s_i$. At $t = 0$, $f_0 = s_j$ has its maximum
at $\bm{x}_j^*$. Since the gradients conflict, $\bm{x}_i^* \neq \bm{x}_j^*$.

For intermediate $t \in (0,1)$, the maximum of $f_t$ traces a path in
$\MC$. By the mountain-pass theorem (a consequence of minimax methods
in critical point theory), there exists at least one saddle point on
$\MC$ separating the basins of attraction of $\bm{x}_i^*$ and
$\bm{x}_j^*$. For the aggregated fitness $f = \Phi(\bm{s})$, each
basin contains a local maximum, hence $f$ has at least two. \qedhere
\end{proof}
\end{theorem}

\TF{3.3}

\begin{definition}[Style]\label{def:3.4}
A \emph{style} is a connected component of the occupied region $B$
concentrated around a local maximum $\bm{x}^*$ of $f$:
\[
  S(\bm{x}^*, \varepsilon, \delta) \;=\;
  \bigl\{\, \bm{x} \in B \;\big|\;
    \norm{\bm{x} - \bm{x}^*} < \varepsilon
    \;\text{and}\;
    f(\bm{x}) > f(\bm{x}^*) - \delta
  \,\bigr\}
\]
for suitable $\varepsilon > 0$, $\delta > 0$. A style is an attraction
basin in the fitness landscape, restricted to realized objects.
Styles are statistical phenomena, not categorical ones: boundaries are
diffuse, and membership is graded. \TF{3.4, 1.2, 1.21}
\end{definition}

\begin{theorem}[Analyst-dependence of style decomposition]\label{thm:3.5}
The decomposition of $B$ into styles depends on the choice of
aggregation function $\Phi$ and weights $\bm{w}$. Different analysts
choosing different $\Phi$ may identify different numbers and boundaries
of styles.

\begin{proof}
By Theorem~\ref{thm:2.8}, different weight vectors $\bm{w}$ produce
different fitness functions $f_{\bm{w}}$ with different local maxima.
Since styles are defined relative to local maxima of $f$
(Definition~\ref{def:3.4}), and $f$ depends on $\Phi$ and $\bm{w}$,
the style decomposition is a function of the analyst's choice.
What is \emph{not} analyst-dependent is the morphospace $\MC$ itself,
the selection pressures $s_i$, and the empirical distribution over $B$.
The landscape is a construction; the pressures are real. \qedhere
\end{proof}
\end{theorem}

\TF{3.41}

\begin{definition}[Path (wandering)]\label{def:3.6}
A \emph{formgiving history} (design trajectory) is a continuous path
in morphospace:
\[
  \gamma : [0, T] \longrightarrow \MC
\]
with $\gamma(0) = \bm{x}^0$ (initial form) and $\gamma(T) = \bm{x}^T$
(final form), such that each step is local:
$\norm{\gamma(t + dt) - \gamma(t)} < \varepsilon$ for all $t$.
The path is \emph{history-dependent}: the accessible region at time $t$
depends on $\gamma\bigl([0, t]\bigr)$. \TF{3.6}
\end{definition}


\newpage
%% ══════════════════════════════════════════════════════════════════
%% CHAPTER 4: DYNAMICS
%% ══════════════════════════════════════════════════════════════════
\section{Dynamics}\label{sec:dynamics}

\begin{postulate}[Dynamic landscape]\label{post:4.1}
Selection pressures are time-dependent:
\[
  s_i(\bm{x},\, t) : \MC \times \R \longrightarrow \R
\]
and therefore the fitness landscape $\FitLand(t)$ is a time-varying
functional:
\[
  f(\bm{x},\, t) = \Phi\bigl(s_1(\bm{x}, t),\; \ldots,\;
  s_k(\bm{x}, t)\bigr).
\]

\medskip\noindent
\textit{Falsification criterion.}
The postulate fails if the fitness landscape for a class $C$ can be
shown to be topologically unchanged over a period in which the observed
distribution of forms changes in a way not attributable to stochastic
drift. \TF{4.1}
\end{postulate}

\begin{theorem}[Punctuated equilibrium]\label{thm:4.2}
Let $\bm{x}(t)$ be a form occupying a local maximum of $f(\cdot, t)$
at time $t$. Then:
\begin{enumerate}[label=(\roman*),nosep]
  \item \textbf{Stability:} While $\partial f / \partial t$ is small
    relative to the curvature $\norm{\nabla^2 f}$ at the peak,
    $\bm{x}(t)$ varies continuously and slowly (the style evolves
    gradually).
  \item \textbf{Catastrophe:} When $\partial^2 f / \partial x^2$ changes
    sign at the peak (the maximum becomes a saddle point), $\bm{x}(t)$
    undergoes a discontinuous jump to a new local maximum (style shift).
\end{enumerate}

\begin{proof}
\textit{(i)}\; By the implicit function theorem, if $\bm{x}^*(t)$ is a
non-degenerate local maximum of $f(\cdot, t)$---meaning the Hessian
$H = \nabla^2 f(\bm{x}^*, t)$ is negative definite---then the map
$t \mapsto \bm{x}^*(t)$ is smooth. The velocity of the peak is
\[
  \frac{d\bm{x}^*}{dt} = -H^{-1} \frac{\partial \grad f}{\partial t}
  \bigg|_{\bm{x}^*}.
\]
While $\partial f / \partial t$ is small and $H$ remains negative
definite (large curvature), $d\bm{x}^*/dt$ is small: the peak drifts
slowly. Forms near this peak exhibit gradual stylistic evolution.

\textit{(ii)}\; When a parameter $t$ reaches a critical value $t_c$ at
which $\det(H) = 0$, the local maximum degenerates. In the generic case
(fold catastrophe in the sense of Thom), two critical points---a maximum
and a saddle---coalesce and annihilate. Beyond $t_c$, the original basin
of attraction disappears, and trajectories that were stable near
$\bm{x}^*(t)$ are redirected to a distant local maximum.

This is the mathematical structure underlying punctuated equilibrium
(\textit{cf.} Gould; Bowler): long periods of gradual drift punctuated
by rapid transitions when the landscape topology changes. \qedhere
\end{proof}
\end{theorem}

\TF{4.11, 4.12, 4.13}

\begin{definition}[Landscape memory]\label{def:4.3}
Every realized form $\bm{y} \in B$ modifies the future selection
pressures via an \emph{impact function} $\eta_i$:
\[
  s_i(\bm{x},\, t + dt) \;=\;
  s_i(\bm{x},\, t) + \eta_i(\bm{y},\, \bm{x})
\]
where $\bm{y} = \gamma(t)$ is the currently realized form and
$\eta_i : \MC \times \MC \to \R$ is continuous with
$\eta_i(\bm{y}, \bm{y}) \geq 0$ (a realized form at least weakly
reinforces its own position).

The landscape \emph{remembers} past forms: current selection pressures
are functions of the entire realized trajectory $\gamma([0, t])$,
not just the current state. This creates co-evolutionary feedback
between form and environment. \TF{4.2, 4.21, 4.22}
\end{definition}


\newpage
%% ══════════════════════════════════════════════════════════════════
%% CHAPTER 5: EMPIRICAL -- MATERIAL SIGNATURES
%% ══════════════════════════════════════════════════════════════════
\section{Empirical: Material Signatures}\label{sec:empirical}

\begin{postulate}[Geometric material signature]\label{post:5.1}
For every material $m$ used in functional class $C$, the conditional
distribution of geometric properties given material is non-uniform:
\[
  P(\bm{x} \mid m) \;\neq\; \mathrm{Uniform}(\MC)
  \qquad \forall\; m.
\]
That is, each material induces a characteristic geometric \emph{profile}
---a region of morphospace that it favours.

\medskip\noindent
\textit{Falsification criterion.}
The postulate fails if the same geometric distribution is observed
regardless of material, within the same functional class and the same
set of other selection pressures. \TF{8.1}
\end{postulate}

\begin{theorem}[Information theorem]\label{thm:5.2}
Within a functional class $C$ with at least two material categories,
material carries strictly more mutual information with realized geometry
than function does:
\[
  \MI(G;\, M) \;>\; \MI(G;\, \mathrm{Func}) \;=\; 0.
\]

\begin{proof}
\textit{Part 1: $\MI(G;\, \mathrm{Func}) = 0$.}\;
By definition of a functional class, all objects in $C$ share the same
function. Therefore $\mathrm{Func}$ is a constant random variable within
$C$: $\Ent(\mathrm{Func} \mid C) = 0$. By the data-processing
inequality and the property that mutual information with a constant is
zero:
\[
  \MI(G;\, \mathrm{Func}) = \Ent(G) - \Ent(G \mid \mathrm{Func})
  = \Ent(G) - \Ent(G) = 0.
\]

\textit{Part 2: $\MI(G;\, M) > 0$.}\;
By Postulate~\ref{post:5.1}, $P(\bm{x} \mid m) \neq P(\bm{x})$ for
at least one material $m$ (the material-conditional distribution differs
from the marginal). This is equivalent to $G$ and $M$ not being
independent. By the fundamental property of mutual information:
\[
  \MI(G;\, M) = 0 \;\iff\; G \perp\!\!\!\perp M.
\]
Since $G \not\perp\!\!\!\perp M$ (by Postulate~\ref{post:5.1}), we
have $\MI(G;\, M) > 0$.

\textit{Combining:}\; $\MI(G;\, M) > 0 = \MI(G;\, \mathrm{Func})$.
\qedhere
\end{proof}
\end{theorem}

\TF{8.11, 8.12, 8.13, 8.14}

\begin{remark}[Empirical values from STOLAR dataset]
In the dataset of $n = 93$ chairs from a Norwegian national museum
(materials: mahogany, steel, laminate, beech; period 1280--2024):
\begin{align*}
  \MI(G;\, M) &= 0.382 \;\text{bits}, \\
  \MI(G;\, \mathrm{Time}) &= 0.168 \;\text{bits}, \\
  \MI(G;\, \mathrm{Geography}) &= 0.079 \;\text{bits}, \\
  \MI(G;\, \mathrm{Func}) &= 0 \;\text{bits (by definition).}
\end{align*}
Material is the strongest predictor. This is the most concentrated
empirical argument against the deterministic axiom ``form follows
function.'' \TF{2.61, 2.62, 8.11--8.18}
\end{remark}

\newpage
%% ══════════════════════════════════════════════════════════════════
%% CHAPTER 6: NAVIGATOR
%% ══════════════════════════════════════════════════════════════════
\section{Navigator}\label{sec:navigator}

\begin{definition}[Navigator]\label{def:6.1}
A \emph{navigator} is a triple
$\Nav = (G,\; \mu,\; \alpha)$
where:
\begin{enumerate}[label=(\alph*),nosep]
  \item $G \subseteq \MC$ is a \emph{goal set} (target region): a
    non-empty subset of morphospace that the navigator steers toward.
  \item $\mu : \MC \to \R_{\geq 0}$ is a \emph{distance function}:
    $\mu(\bm{x}) = d(\bm{x},\, G)$ for a metric $d$ on $\MC$, with
    $\mu(\bm{x}) = 0 \iff \bm{x} \in G$.
  \item $\alpha : \MC \to T\MC$ is a \emph{steering field}: a vector
    field on $\MC$ satisfying
    \[
      \inner{\alpha(\bm{x})}{\!-\grad\mu(\bm{x})} > 0
      \qquad \text{whenever } \mu(\bm{x}) > 0.
    \]
    That is, $\alpha$ always has a component pointing toward decreasing
    distance to $G$.
\end{enumerate}

\medskip\noindent
Three conditions are required: (a)~the system has a goal state,
(b)~it can register distance from the goal, and (c)~it can adjust
toward the goal. No consciousness, intention, or nervous system is
required. A thermostat, a gradient-descent algorithm, and a chemical
reaction driven toward equilibrium are all navigators. \TF{5.1, 5.11}
\end{definition}

\begin{definition}[Cognitive light cone]\label{def:6.2}
The \emph{cognitive light cone} (boundary) of a navigator $\Nav$ at
time $t$ is the maximal region of morphospace--time in which $\Nav$ can
represent states and steer:
\[
  \partial\Nav \;=\;
  \bigl\{\, (r,\, \tau) \in \MC \times \R
    \;\big|\;
    \Nav \text{ can represent and work toward states in } (r, \tau)
  \,\bigr\}.
\]
The spatial extent determines what problems $\Nav$ can address; the
temporal extent determines its planning horizon. \TF{5.2, 5.21}
\end{definition}

\begin{definition}[Intelligence]\label{def:6.3}
The \emph{intelligence} of a navigator $\Nav$ is its capacity to avoid
local optima in the fitness landscape:
\[
  I(\Nav) \;\propto\;
  \max\bigl\{\, d\bigl(\gamma_\Nav(t),\; \bm{x}_{\mathrm{local}}\bigr)
    \;\big|\;
    \gamma_\Nav \text{ converges to } \bm{x}_{\mathrm{global}}
  \,\bigr\}
\]
where $\bm{x}_{\mathrm{local}}$ is a local optimum and
$\bm{x}_{\mathrm{global}}$ is a better alternative.
Intelligence is the ability to follow detours---to accept temporarily
lower fitness in order to reach a higher peak. It is scale-invariant:
the definition applies equally to an enzyme avoiding a conformational
trap and to a civilization investing in infrastructure. \TF{5.3, 5.31}
\end{definition}


\newpage
%% ══════════════════════════════════════════════════════════════════
%% CHAPTER 7: SUBSTRATE INDEPENDENCE
%% ══════════════════════════════════════════════════════════════════
\section{Substrate Independence}\label{sec:substrate}

\begin{postulate}[Substrate-independent navigation]\label{post:7.1}
The conditions (a), (b), (c) in Definition~\ref{def:6.1} can be
satisfied by any physical system---biological, algorithmic, mechanical,
or distributed---that:
\begin{enumerate}[label=(\roman*),nosep]
  \item has states (positions in some configuration space),
  \item can register distance from a target configuration, and
  \item can adjust its state in response.
\end{enumerate}
Navigation is a property of the system's \emph{functional organization},
not of the material it is made from.

\medskip\noindent
\textit{Falsification criterion.}
The postulate fails if it can be shown that at least one of conditions
(a)--(c) necessarily requires a specific physical substrate (e.g.,
organic matter)---that is, if there exist formgiving processes where a
human navigator consistently produces forms that no algorithmic navigator
can produce, \emph{not} because of computational limits or sensory
resolution, but because navigation itself requires organic substrate.
\TF{5.12}
\end{postulate}

\begin{theorem}[Substrate independence of morphospace position]%
\label{thm:7.2}
The position $\bm{x} \in \MC$ occupied by a realized object is
determined by its measurable properties, not by the process that
produced it. Two objects with identical attribute vectors occupy the
same point in $\MC$ regardless of whether they were made by hand, by
machine, by algorithm, or by biological growth.

\begin{proof}
By Definition~\ref{def:1.1}, the attribute vector $\bm{x}(a)$ depends
only on the measurable geometric properties $x_1(a), \ldots, x_n(a)$.
These are intrinsic to the object $a$, not to its production history.
If two objects $a$ (handmade) and $b$ (machine-made) have
$x_i(a) = x_i(b)$ for all $i$, then $\bm{x}(a) = \bm{x}(b)$,
hence they occupy the same point.

This does \emph{not} mean the production process is irrelevant: it
determines which region of $\MC$ is accessible (the substrate constrains
the trajectory, \S\ref{sec:variational}). But the position itself is
substrate-invariant. \qedhere
\end{proof}
\end{theorem}

\TF{1.4, 1.41, 5.12}

\newpage
%% ══════════════════════════════════════════════════════════════════
%% CHAPTER 8: EMERGENCE
%% ══════════════════════════════════════════════════════════════════
\section{Emergence}\label{sec:emergence}

\begin{theorem}[Emergent navigation]\label{thm:8.1}
Let $\Nav^1, \Nav^2, \ldots, \Nav^S$ be a collection of sub-navigators,
each with its own triple $(G_j, \mu_j, \alpha_j)$. The collective system
$\Nav = \{\Nav^1, \ldots, \Nav^S\}$ can realize a navigator triple
$(G, \mu, \alpha)$ where $G$, $\mu$, and $\alpha$ are \emph{not}
functions of any single sub-navigator's triple, but emerge from the
interaction dynamics.

\begin{proof}
Define the collective state as $\bm{X} = (\bm{x}^1, \ldots, \bm{x}^S)
\in {(\MC)}^S$. Each sub-navigator $\Nav^j$ generates a flow
$\dot{\bm{x}}^j = \alpha_j(\bm{x}^j)$ in its own morphospace.
When sub-navigators interact---e.g., through shared selection pressures,
communication, or physical coupling---the collective dynamics become:
\[
  \dot{\bm{x}}^j = \alpha_j(\bm{x}^j)
  + \sum_{l \neq j} \phi_{jl}(\bm{x}^j, \bm{x}^l)
\]
where $\phi_{jl}$ encodes the interaction between sub-navigators $j$
and $l$.

The collective system exhibits a macroscopic goal $G$, distance $\mu$,
and steering $\alpha$ if there exists a coarse-graining map
$\Pi : {(\MC)}^S \to \MC$ such that the projected dynamics
$\dot{\bm{y}} = \Pi_* \dot{\bm{X}}$ satisfy Definition~\ref{def:6.1}
with $(G, \mu, \alpha)$ that cannot be decomposed as functions of any
single $(G_j, \mu_j, \alpha_j)$.

\textit{Existence.}\;
Xenobots---synthetic organisms assembled from frog embryonic cells
without evolutionary history---navigate toward coherent morphological
configurations. The navigation competency is not encoded in any single
cell's genome but emerges from intercellular bioelectric signaling
dynamics (Murugan et al., 2022; Blackiston et al., 2015). This
demonstrates that the collective triple $(G, \mu, \alpha)$ is a real
emergent property, not reducible to individual cell navigation.

In the design context: the final form of a collaboratively designed
object was not specified by any single participant but emerged from the
negotiation between them. \qedhere
\end{proof}
\end{theorem}

\TF{5.6, 5.61, 5.62}


\newpage
%% ══════════════════════════════════════════════════════════════════
%% CHAPTER 9: VARIATIONAL PRINCIPLES
%% ══════════════════════════════════════════════════════════════════
\section{Variational Principles}\label{sec:variational}

\begin{postulate}[Path-dependent fitness]\label{post:9.1}
The fitness value of a realized form depends not only on its position
$\bm{x} \in \MC$ but on the entire path $\gamma : [0, T] \to \MC$
taken to reach it. Define the \emph{path fitness functional}:
\[
  S[\gamma] \;=\; \int_0^T
    g\!\left(\gamma(t),\; \dot{\gamma}(t),\; t\right) dt
\]
where $g : \MC \times T\MC \times \R \to \R$ is a Lagrangian that
depends on the position $\gamma(t)$, velocity $\dot{\gamma}(t)$
(rate of change), and time $t$.

\medskip\noindent
\textit{Interpretation.}\;
Two identical chairs occupying the same point in morphospace but with
different production histories (e.g., original vs.\ copy) have
different cultural value, different market position, and different
potential to influence future designs. Position alone is insufficient;
the path matters.

\medskip\noindent
\textit{Falsification criterion.}\;
The postulate fails if identical forms with different production
histories consistently have the same survival and reproduction rates
in the design ecology. \TF{7.1}
\end{postulate}

\begin{theorem}[Stationary path]\label{thm:9.2}
If realized design trajectories are stationary points of $S[\gamma]$
(the conjecture $\delta S[\gamma] = 0$), then they satisfy the
Euler--Lagrange equation:
\[
  \frac{\partial g}{\partial \gamma}
  - \frac{d}{dt}\frac{\partial g}{\partial \dot{\gamma}} = 0.
\]

\begin{proof}
This is a direct application of the calculus of variations. Let
$\gamma_\varepsilon(t) = \gamma(t) + \varepsilon \, \eta(t)$ be a
variation with $\eta(0) = \eta(T) = 0$. The stationarity condition is:
\[
  \frac{d}{d\varepsilon}\bigg|_{\varepsilon=0} S[\gamma_\varepsilon]
  = 0.
\]
Computing:
\begin{align*}
  \frac{d}{d\varepsilon}\bigg|_{0} S[\gamma_\varepsilon]
  &= \int_0^T \left(
    \frac{\partial g}{\partial \gamma} \cdot \eta
    + \frac{\partial g}{\partial \dot{\gamma}} \cdot \dot{\eta}
  \right) dt \\
  &= \int_0^T \left(
    \frac{\partial g}{\partial \gamma}
    - \frac{d}{dt}\frac{\partial g}{\partial \dot{\gamma}}
  \right) \cdot \eta \; dt
  + \left[\frac{\partial g}{\partial \dot{\gamma}} \cdot \eta
  \right]_0^T.
\end{align*}
The boundary term vanishes since $\eta(0) = \eta(T) = 0$. Since $\eta$
is arbitrary, the fundamental lemma of the calculus of variations
gives the Euler--Lagrange equation. \qedhere
\end{proof}
\end{theorem}

\TF{7.11}

\begin{theorem}[Path modifies landscape]\label{thm:9.3}
The path-dependent fitness and landscape memory (Definition~\ref{def:4.3})
imply that giving form is not finding the best position but finding the
best \emph{path}: the process modifies what outcomes are possible.

\begin{proof}
By Definition~\ref{def:4.3}, $s_i(\bm{x}, t+dt)$ depends on the
realized form $\bm{y} = \gamma(t)$ via the impact function $\eta_i$.
The fitness landscape at time $T$ is therefore a functional of the
entire trajectory:
\[
  f(\bm{x}, T) = \Phi\!\left(
    s_1\!\left(\bm{x},\, 0\right) +
    \int_0^T \eta_1(\gamma(\tau), \bm{x})\, d\tau,\;
    \ldots
  \right).
\]
The final landscape---and therefore the set of available optima---depends
on the path taken. Two different trajectories $\gamma$ and $\gamma'$
arriving at the same final point $\bm{x}^T$ may face different future
landscapes, different available optima, and different survival
prospects. The process is not merely a means to the result; the process
\emph{modifies} what results are possible. \qedhere
\end{proof}
\end{theorem}

\TF{7.12}

\begin{theorem}[Spreading relation]\label{thm:9.4}
Let $\sigma^2[\gamma]$ denote the variance of the form distribution
around a given attractor in $\MC$, and let $\norm{\bm{s}}$ denote the
aggregate strength of selection pressures at that attractor. Then:
\[
  \sigma^2 \;\propto\; \frac{1}{\norm{\bm{s}}}.
\]
Stronger selection produces narrower distributions (tighter styles);
weaker selection produces broader distributions (more variation).

\begin{proof}
Model the stationary distribution of forms near a local maximum
$\bm{x}^*$ as a Boltzmann-like distribution:
\[
  P(\bm{x}) \;\propto\; \exp\!\bigl(\beta \, f(\bm{x})\bigr)
\]
where $\beta > 0$ controls the selection intensity (analogous to
inverse temperature). Near the maximum, expand $f$ to second order:
\[
  f(\bm{x}) \approx f(\bm{x}^*) - \tfrac{1}{2}
  (\bm{x} - \bm{x}^*)^\top H (\bm{x} - \bm{x}^*)
\]
where $H = -\nabla^2 f(\bm{x}^*)$ is positive definite (since
$\bm{x}^*$ is a maximum). Then:
\[
  P(\bm{x}) \propto \exp\!\left(
    -\tfrac{\beta}{2}
    (\bm{x} - \bm{x}^*)^\top H (\bm{x} - \bm{x}^*)
  \right).
\]
This is a multivariate Gaussian with covariance
$\Sigma = (\beta H)^{-1}$. The variance along any axis is
$\sigma_i^2 = (\beta \, \lambda_i)^{-1}$ where $\lambda_i$ are the
eigenvalues of $H$. Since $H$ reflects the curvature of the aggregated
selection pressures, $\lambda_i \propto \norm{\bm{s}}$, giving
$\sigma^2 \propto 1/\norm{\bm{s}}$. \qedhere
\end{proof}
\end{theorem}

\TF{7.2, 7.21}

\begin{definition}[Substrate-dependent distribution]\label{def:9.5}
The probability of form $\bm{x}$ given substrate (material, tool,
process) is:
\[
  P(\bm{x} \mid \text{substrate}) \;\neq\;
  P(\bm{x} \mid \text{substrate}')
\]
for distinct substrates. Each substrate---each tradition, each tool,
each material---generates a different distribution over $\MC$. A
substrate shift changes not just \emph{how} forms are made but
\emph{which} forms are discovered. \TF{7.3}
\end{definition}


\newpage
%% ══════════════════════════════════════════════════════════════════
%% CHAPTER 10: CONSEQUENCES
%% ══════════════════════════════════════════════════════════════════
\section{Consequences}\label{sec:consequences}

\begin{theorem}[Provisional compromise]\label{thm:10.1}
Every realized form is a \emph{provisional} compromise: optimal for the
selection pressures at the time of its creation, but generically
suboptimal for any later configuration of pressures.

\begin{proof}
By Theorem~\ref{thm:2.7}, every realized form $\bm{x} \in B$ is a
compromise: it balances conflicting selection pressures but does not
maximize all simultaneously. By Postulate~\ref{post:4.1}, the pressures
$s_i(\bm{x}, t)$ change over time.

Let $\bm{x}^*(\tau)$ be the local maximum of
$f(\cdot, \tau)$ nearest to $\bm{x}$ at time $\tau$ of creation. At a
later time $t > \tau$, the pressures have shifted:
$\bm{s}(\cdot, t) \neq \bm{s}(\cdot, \tau)$. The point
$\bm{x}^*(\tau)$ is generically \emph{not} a critical point of
$f(\cdot, t)$, since
\[
  \grad f(\bm{x}^*(\tau),\, t) =
  \grad f(\bm{x}^*(\tau),\, \tau) +
  \int_\tau^t \frac{\partial}{\partial t'} \grad f(\bm{x}^*(\tau),\, t')\, dt'
  \neq \bm{0}
\]
unless the integral vanishes identically (a non-generic condition).
Therefore the form that was optimal at $\tau$ is no longer optimal at $t$.
No form is final; every object is a snapshot of a dynamic system.
\qedhere
\end{proof}
\end{theorem}

\TF{7.4}

\begin{theorem}[Durable traditions]\label{thm:10.2}
Traditions with more independent adaptive mechanisms (more selection
pressures addressed by competent sub-navigators) are more robust to
landscape changes.

\begin{proof}
Let a tradition $\mathcal{T}$ be a population of navigators addressing
$k_\mathcal{T}$ distinct selection pressures through independent adaptive
mechanisms. When the landscape shifts (Postulate~\ref{post:4.1}),
typically only a subset of pressures change simultaneously.

A tradition with $k_\mathcal{T}$ independent mechanisms can respond to
a shift in pressure $s_i$ without disrupting its response to pressures
$s_j$ ($j \neq i$), because the mechanisms are independent. Its
adaptation requires adjusting only $\Delta k \leq k_\mathcal{T}$
mechanisms.

Compare traditions $\mathcal{T}_1$ with $k_1$ mechanisms and
$\mathcal{T}_2$ with $k_2 > k_1$ mechanisms. Under a landscape shift
affecting one pressure, $\mathcal{T}_2$ adjusts only
$\Delta k / k_2 < \Delta k / k_1$ of its total adaptive capacity.
The larger repertoire provides redundancy: a craft tradition that
combines material understanding, cultural sensitivity, and individual
innovation adapts better than one relying on a single mechanism.
\qedhere
\end{proof}
\end{theorem}

\TF{7.5}

\begin{theorem}[Falsifiability of the framework]\label{thm:10.3}
The framework of Tractatus Formae is falsifiable: if a single selection
pressure $s_1$ is sufficient to explain all observed form variation
within a functional class $C$, then Postulate~\ref{post:2.3} fails, and
with it Theorems~\ref{thm:2.4}, \ref{thm:2.7}, \ref{thm:2.8},
\ref{thm:3.3}, \ref{thm:3.5}, and \ref{thm:10.1}.

\begin{proof}
The deductive chain is:
\begin{enumerate}[nosep]
  \item Postulate~\ref{post:2.3} (multiple independent pressures)
    $\implies$ Theorem~\ref{thm:2.4} (variation not explained by
    function).
  \item Postulates~\ref{post:2.3} + \ref{post:2.6} (conflicting
    gradients) $\implies$ Theorem~\ref{thm:2.7} (compromise) $\implies$
    Theorem~\ref{thm:2.8} (multiple compromises).
  \item Theorems~\ref{thm:2.7} + \ref{thm:2.8} $\implies$
    Theorem~\ref{thm:3.3} (multiple peaks) $\implies$
    Definition~\ref{def:3.4} (styles exist).
  \item Theorem~\ref{thm:3.3} + Definition~\ref{def:3.4} $\implies$
    Theorem~\ref{thm:3.5} (analyst dependence of style decomposition).
  \item Postulate~\ref{post:4.1} (dynamic landscape) + Theorem~\ref{thm:2.7}
    $\implies$ Theorem~\ref{thm:10.1} (provisional compromise).
\end{enumerate}

If any of the five postulates can be empirically disproven, the
corresponding branch of the deductive tree collapses. The framework
is valid precisely because it can fail. The deterministic axiom
``form follows function'' was unfalsifiable (Michl) and therefore
not a scientific claim. This framework makes quantitative predictions
(Theorem~\ref{thm:9.4}, Postulate~\ref{post:5.1}) that can be tested
against data.

That the text \emph{can} be felled is the reason it is valid.
\qedhere
\end{proof}
\end{theorem}

\TF{7.7}


\newpage
%% ══════════════════════════════════════════════════════════════════
%% APPENDIX: CORRESPONDENCE TABLE
%% ══════════════════════════════════════════════════════════════════
\section*{Appendix: Correspondence Table}
\addcontentsline{toc}{section}{Appendix: Correspondence Table}

The following table maps every formal element in this document to its
counterpart in \textit{Tractatus Formae}.

\medskip
\begin{longtable}{@{}lll@{}}
\toprule
\textsc{Mathematica} & \textsc{Type} & \textsc{Tractatus Formae} \\
\midrule
\endfirsthead
\toprule
\textsc{Mathematica} & \textsc{Type} & \textsc{Tractatus Formae} \\
\midrule
\endhead
\bottomrule
\endfoot

Def.~1.1 (Morphospace) & Definition & 1.1 \\
Def.~1.2 (Constraints) & Definition & 1.1 \\
Def.~1.3 (Region typology) & Definition & 1.13 \\
Prop.~1.4 (Explanation requirement) & Proposition & 1.14 \\
\addlinespace
Def.~2.1 (Selection pressure) & Definition & 2.1 \\
Def.~2.2 (Pressure vector) & Definition & 2.14 \\
Post.~2.3 (Multiple pressures) & Postulate & 2.2 \\
Thm.~2.4 (Variation under function) & Theorem & 2.21 \\
Def.~2.5 (Functional delimitation) & Definition & 2.4 \\
Post.~2.6 (Conflicting gradients) & Postulate & 2.3 \\
Thm.~2.7 (Compromise) & Theorem & 2.31 \\
Thm.~2.8 (Multiple compromises) & Theorem & 2.32 \\
\addlinespace
Def.~3.1 (Fitness function) & Definition & 3.1 \\
Def.~3.2 (Fitness landscape) & Definition & 3.2 \\
Thm.~3.3 (Multiple peaks) & Theorem & 3.3 \\
Def.~3.4 (Style) & Definition & 3.4 \\
Thm.~3.5 (Analyst dependence) & Theorem & 3.41 \\
Def.~3.6 (Path / wandering) & Definition & 3.6 \\
\addlinespace
Post.~4.1 (Dynamic landscape) & Postulate & 4.1 \\
Thm.~4.2 (Punctuated equilibrium) & Theorem & 4.11 \\
Def.~4.3 (Landscape memory) & Definition & 4.2 \\
\addlinespace
Post.~5.1 (Geometric signature) & Postulate & 8.1 \\
Thm.~5.2 (Information theorem) & Theorem & 8.11 \\
\addlinespace
Def.~6.1 (Navigator) & Definition & 5.1 \\
Def.~6.2 (Cognitive light cone) & Definition & 5.2 \\
Def.~6.3 (Intelligence) & Definition & 5.3 \\
\addlinespace
Post.~7.1 (Substrate independence) & Postulate & 5.12 \\
Thm.~7.2 (Position invariance) & Theorem & 1.4, 5.12 \\
\addlinespace
Thm.~8.1 (Emergent navigation) & Theorem & 5.62 \\
\addlinespace
Post.~9.1 (Path-dependent fitness) & Postulate & 7.1 \\
Thm.~9.2 (Stationary path) & Theorem & 7.11 \\
Thm.~9.3 (Path modifies landscape) & Theorem & 7.12 \\
Thm.~9.4 (Spreading relation) & Theorem & 7.2 \\
Def.~9.5 (Substrate-dependent distribution) & Definition & 7.3 \\
\addlinespace
Thm.~10.1 (Provisional compromise) & Theorem & 7.4 \\
Thm.~10.2 (Durable traditions) & Theorem & 7.5 \\
Thm.~10.3 (Falsifiability) & Theorem & 7.7 \\
\end{longtable}


\newpage
%% ══════════════════════════════════════════════════════════════════
%% APPENDIX: DEPENDENCY GRAPH
%% ══════════════════════════════════════════════════════════════════
\section*{Appendix: Logical Dependency Graph}
\addcontentsline{toc}{section}{Appendix: Logical Dependency Graph}

The deductive structure of the theory, showing which results depend on
which postulates and earlier theorems:

\medskip
\begin{small}
\begin{verbatim}
                    +-- Def 1.1 (Morphospace)
                    |     +-- Def 1.2 (Constraints)
                    |         +-- Def 1.3 (Regions)
                    |             +-- Prop 1.4 (Explanation)
                    |
   POSTULATE 2.3 --+-- Thm 2.4 (Variation /= function)
   (k >= 2)        |
                    |          +-- Thm 2.7 (Compromise)
   POSTULATE 2.6 --+----------+
   (grad conflict)  |          +-- Thm 2.8 (Multiple compromises)
                    |                |
                    |          Thm 3.3 (Multiple peaks)
                    |           +-- Def 3.4 (Style)
                    |           +-- Thm 3.5 (Analyst dependence)
                    |
   POSTULATE 4.1 --+-- Thm 4.2 (Punctuated equilibrium)
   (dynamic)        |   +-- Def 4.3 (Landscape memory)
                    |
                    +-- Thm 10.1 (Provisional compromise)
                    |       [from 2.7 + 4.1]
                    |
                    +-- Thm 10.2 (Durable traditions)
                            [from 10.1 + 6.1]

   POSTULATE 5.1 ---- Thm 5.2 (Information theorem)
   (material sig.)

   POSTULATE 7.1 ---- Thm 7.2 (Position invariance)
   (substrate-ind.)    +-- Thm 8.1 (Emergent navigation)

   POSTULATE 9.1 --+-- Thm 9.2 (Stationary path / E-L)
   (path-dep.)      +-- Thm 9.3 (Path modifies landscape)
                    +-- Thm 9.4 (Spreading relation)

   Thm 10.3 (Falsifiability)
     [meta-thm: if Post 2.3 fails -> 2.4, 2.7, 2.8, 3.3, 10.1 fall]
\end{verbatim}
\end{small}

\medskip\noindent
\textbf{Summary of formal inventory:}
\begin{itemize}[nosep]
  \item 14 definitions (Def.~1.1--9.5)
  \item 7 postulates (Post.~2.3, 2.6, 4.1, 5.1, 7.1, 9.1 + the
    explanation principle 1.4)
  \item 13 theorems (Thm.~2.4--10.3)
  \item 1 meta-theorem on falsifiability (Thm.~10.3)
\end{itemize}

\vfill
\begin{center}
\rule{4cm}{0.4pt}

\medskip
{\small\itshape This document is the formal companion to}\\
{\small\scshape Tractatus Formae}\\
{\small\itshape Iver Raknes Finne, 2026}
\end{center}

%% ── Bibliography ─────────────────────────────────────────────────
\newpage
\printbibliography[title={References}]

\end{document}
